\documentclass[11pt,letterpaper]{article}

\usepackage{latexsym}
\usepackage{amssymb}
\usepackage{amsmath}
\usepackage{graphicx}
\usepackage[pdfpagemode=None,pdfstartview=FitH,pdfview=FitH,colorlinks=true]%
 {hyperref}

\newtheorem{theorem}{Theorem}[section]
\newcommand{\qi}{\ensuremath{\mathit{qi}}}

\pagestyle{headings}
\bibliographystyle{alpha}

\title{Theora I Specification}
\author{Xiph.org Foundation}
\date{\today}

\begin{document}

\maketitle
\tableofcontents
\newpage

\section{Introduction and Description}

This section provides a high level description of the Theora codec's
 construction.
A bit-by-bit specification appears beginning in Section~\ref{sec:bitpacking}.
The later sections assume a high-level understanding of the Theora decode
 process, which is provided below.

\subsection{Overview}

Theora is a general purpose, lossy video codec.
It is based on the VP3 video codec produced by On2 Technologies
 (\url{http://www.on2.com/}).
On2 Technologies donated the VP3.2 source code to the Xiph.org 
Foundation and it was released under a BSD license. On2 also made an 
irrevocable, royalty-free license grant for any patent claims it might 
have over the software and any derivatives.
No formal specification exists for the VP3 format beyond this source code,
 though Mike Melanson maintains a detailed description \cite{Mel04}.
Portions of this specification were adopted from his text with permission.

\subsubsection{VP3 and Theora}

Theora contains a superset of the features that were available in the original
 VP3 codec.
Content encoded with VP3.2 can be losslessly transcoded into the Theora format.
%TODO: what about VP3.1 etc? source tables all say 'VP31'
Theora content cannot, in general, be losslessly transcoded into the VP3
 format.
If a feature is not available in the original VP3 format, this is
 mentioned when that feature is defined.
A complete list of these features appears in Appendix~REF.

\subsubsection{Video Formats}

Theora I currently supports progressive video data of arbitrary dimensions in
 one of several $Y'C_bC_r$ color spaces.
The precise definition the color spaces supported appears in Section~REF.
Three different chroma subsampling formats are supported: 4:2:0, 4:2:2,
 and 4:4:4.
The precise details of each of these formats and their sampling locations are
 described in Section~REF.

The Theora I format does not support interlaced material, bit-depths larger
 than 8 bits per component, nor alternate color spaces such as RGB or
 arbitrary multi-channel spaces. Black and white content can be 
efficiently encoded because the uniform chromaticity planes compress 
well.
Support for interlaced material is planned for a future version.
Support for increased bit depths or additional color spaces is not 
planned.

\subsubsection{Classification}

Theora I is a block-based lossy transform codec that utilizes an
 $8\times 8$ Type-II Discrete Cosine Transform and block-based motion
 compensation.
This places it in the same class of codecs as MPEG-1, -2, -4, and H.263.
The details of how individual blocks are organized and how DCT coefficients are
 organized in the bitstream differ stubstantially from these codecs, however.
Theora supports only intra frames (I frames in MPEG) and inter frames (P frames
 in MPEG).
There is no equivalent to the bi-predictive frames (B frames) 
found in MPEG codecs.

\subsubsection{Assumptions}

The Theora codec design assumes a complex, psychovisually-aware encoder and a
 simple, low-complexity decoder.
%TODO: Talk more about implementation complexity.

Theora provides none of its own framing, synchronization, or protection against
 transmission errors; it is solely a method of accepting input video 
frames and compressing
 these frames into raw, unformatted `packets'.
The decoder then accepts these raw packets in sequence, decodes them,
 and synthesizes a fascimile of the original video frames.
Theora is a free-form variable bit rate (VBR) codec, and packets have no
 minimum size, maximum size, or fixed/expected size.

Theora packets are thus intended to be used with a transport mechanism that
 provides free-form framing, synchronization, positioning, and error correction
 in accordance with these design assumptions, such as Ogg (for file transport)
 or RTP (for network multicast).
For the purposes of a few examples in this document, we will assume that Theora
 is embedded in an Ogg stream specifically, although this is by no means a
 requirement or fundamental assumption in the Theora design.

The specification for embedding Theora into an Ogg transport stream is 
given in
 Appendix~REF.

\subsubsection{Codec Setup and Probability Model}

Theora's heritage is the proprietary commerical codec VP3, and it retains
 a fair amount of inflexibility when compared to Vorbis, the first 
Xiph.org codec.
However, to provide additional scope for encoder improvement,
Theora adopts some of the configurable aspects of decoder setup that 
are present in Vorbis.

Theora makes the same controversial design decision that Vorbis made to include
 the entire probability model for the DCT coefficients and all the quantization
 parameters in the bitstream headers.
This is often several hundred fields.
This makes it impossible to begin decoding at any frame in the stream without
 having previously fetched the codec info and codec setup headers.

\begin{verse}
{\bf Note:} Theora {\em can} initiate decode at an arbitrary intra-frame packet
 within a bitstream so long as the codec has been initialized with the setup
 headers.
\end{verse}

Thus, Theora headers are both required for decode to begin and relatively large
 as bitstream headers go.
The header size is unbounded, although as a rule-of-thumb of less than 16kB 
 is recommended, and Xiph.org's reference encoder follows this 
 suggestion.
%TODO: Is 8kB enough? My setup header is 7.4kB, that doesn't leave much room
% for comments.
%RG: the lesson from vorbis is that as small as possible is really 
% important in some applications. Practically, what's acceptable
% depends a great deal on the target bitrate. I'd leave 16 kB in the
% spec for now. fwiw more than 1k of comments is quite unusual.

Our own design work indicates that the primary liability of the required header
 is in mindshare; it is an unusual design and thus causes some amount of
 complaint among engineers as this runs against current design trends, and
 points out limitations in some existing software/interface designs.
However, we find that it does not fundamentally limit Theora's suitable
 application space.

\subsubsection{Format Specification}

The Theora format is well-defined by its decode specification; any encoder that
 produces packets that are correctly decoded by an implementation following
 this specification may be considered a proper Theora encoder.
A decoder must faithfully and completely implement the specification defined
 herein %, except where noted,
 to be considered a proper Theora decoder.
Where appropriate, a non-normative description of encoder processes are
 included.
These sections will be marked as such, and a proper Theora encoder is not
 bound to follow them.

%TODO: \subsubsection{Hardware Profile}

\subsection{Decoder Configuration}

Decoder setup consists of configuration of the quantization matrices and the
 Huffman codebooks for the DCT coefficients.
The remainder of the decoding pipeline is not configurable.

\subsubsection{Global Configuration}

The global codec configuration consists of a few video related fields, such as
 frame rate, frame size, picture size and offset, aspect ratio, color space,
 pixel format, and a version number.
The version number is divided into a major version, a minor version, amd a
 minor revision number.
For the format defined in this specification, these are `3', `2', and 
 `0', respectively, in reference to Theora's origin as
 a successor to the VP3.1 format.

\subsubsection{Quantization Matrices}

Theora allows up to 384 different quantization matrices to be defined, one for
 each {\em quantization type} (intra or inter), {\em color plane} ($Y'$, $C_b$,
 or $C_r$), and {\em quantization index}, \qi, which ranges from zero to 63,
 inclusive.
The quantization index generally represents a progressive range of quality
 levels, from low quality near zero to high quality near 63.
However, the interpretation is arbitrary, and it is possible, for example, to
 partition the scale into two completely separate ranges with 32 levels each
 that are meant to represent different classes of source material.

Each quantization matrix is an $8\times 8$ matrix of 16-bit values, which is
 used to quantize the output of the $8\times 8$ DCT.
Quantization matrices are specified using three components: a {\em base matrix}
 and two {\em scale values}.
The first scale value is the {\em DC scale}, which is applied to the DC
 component of the base matrix.
The second scale value is the {\em AC scale}, which is applied to all the other
 components of the base matrix.

There are 64 DC scale values and 64 AC scale values, one for each \qi value.
There is a set of base matrices for each quantization type and each color
 plane.
The bitstream specifies this set by defining a base matrix for a sparse subset
 of the possible \qi values, including at least zero and 63.
The base matrices for the remainder of the \qi values are computed using linear
 interpolation.
This configuration allows the quantization matrices to approximate the complex,
 non-linear processes of the human visual system as the \qi value varies.

Finally, because the in-loop deblocking filter strength depends on the strength
 of the quantization matrices defined in this header, a table of 64 {\em loop
 filter limit values} is defined, one for each \qi value.

The precise specification of how all of this information is decoded appears in
 Section~REF.

\subsubsection{Huffman Codebooks}

Theora uses 80 configurable binary Huffman codes to represent the 32 tokens
 used to encode DCT coefficients.
Each of the 32 token values has a different semantic meaning and is used to
 represent single coefficient values, zero runs, combinations of the two, and
 `end-of-block' markers.

The 80 codes are divided up into five groups of 16, with each group
 corresponding to a set of DCT coefficient indices.
The first group corresponds to the DC coefficient, while the remaining groups
 correspond to different subsets of the AC coefficients.
Within each frame, two 4-bit codebook indices are stored.
The first selects which codebook to use from the DC coefficient group, while
 the second selects which codebook to use from {\em all} of the AC coefficient
 groups.

The precise specification of how the codebooks are decoded appears in
 Section~REF.

\section{Bitpacking Convention}
\label{sec:bitpacking}

\subsection{Overview}

The Theora codec uses relatively unstructured raw packets containing
 binary integer fields of arbitrary width.
Logically, each packet is a bitstream in which bits are written one-by-one by
 the encoder and then read one-by-one in the same order by the decoder.
Most current binary storage arrangements group bits into a native storage unit
 of eight bits (octets), sixteen bits, thirty-two bits, or less commonly other
 fixed sizes.
The Theora bitpacking convention specifies the correct mapping of the logical
 packet bitstream into an actual representation in fixed-width units.

\subsubsection{Octets and Bytes}

In most contemporary architectures, a `byte' is synonymous with an `octect',
 that is, eight bits.
For purposes of the bitpacking convention, a byte implies the smallest native
 integer storage representation offered by a platform.
Modern file systems invariably offer bytes as the fundamental atom of storage.

The most ubiquitous architectures today consider a `byte' to be an octet.
Note, however, that the Theora bitpacking convention is still well defined for
 any native byte size;
an implementation can use the native bit-width of a given storage 
system.
This document assumes that a byte is one octet for purposes of example only.

\subsubsection{Words and Byte Order}

A `word' is an integer size that is a grouped multiple of the byte size.
Most architectures consider a word to be a group of two, four, or eight bytes.
Each byte in the word can be ranked by order of `significance', e.g. the
 significance of the bits in each byte when storing a binary integer in the
 word.
Several byte orderings are possible in a word.
The common ones are
\begin{itemize}
\item{Big-endian:}
in which the most significant byte comes first, e.g. 3-2-1-0,
\item{Little-endian:}
in which the least significant byte comes first, e.g. 0-1-2-3, and less
 commonly
\item{Mixed-endian:}
e.g. 3-1-2-0 or 0-2-1-3.
\end{itemize}

The Theora bitpacking convention specifies storage and bitstream manipulation
 at the byte, not word, level.
Thus host word ordering is of a concern only during optimization, when writing
 code that operates on a word of storage at a time rather than a byte.
Logically, bytes are always encoded and decoded in order from byte zero through
 byte $n$.

\subsubsection{Bit Order}

A byte has a well-defined `least significant' bit (LSb), which is the only bit
 set when the byte is storing the two's complement integer value $+1$.
A byte's `most significant' bit (MSb) is at the opposite end of the byte.
Bits in a byte are numbered from zero at the LSb to $n$ for the MSb, where
 $n=7$ in an octet.

\subsection{Coding Bits into Bytes}

The Theora codec needs to encode arbitrary bit-width integers from zero to 32
 bits wide into packets.
These integer fields are not aligned to the boundaries of the byte
 representation; the next field is read at the bit position immediately
 after the end of the previous field.

The decoder logically unpacks integers by first reading the MSb of a binary
 integer from the logical bitstream, followed by the next most significant
 bit, etc., until the required number of bits have been read.
When unpacking the bytes into bits, the decoder begins by reading the MSb of
 the integer to be read from the most significant unread bit position of the
 source byte, followed by the next-most significant bit position of the
 destination integer, and so on up to the requested number of bits.
Note that this differs from the Vorbis I codec, which
 begins decoding with the LSb of the source integer, reading it from the
 LSb of the source byte.
When all the bits of the current source byte are read, decoding continues with
 the MSb of the next byte.
Any unfilled bits in the last byte of the packet must be cleared to zero by the
 encoder.

\subsubsection{Signedness}

The binary integers decoded by the above process may be either signed or
 unsigned.
This varies from integer to integer, and this specification 
 indicates how each value should be interpreted as it is read.
That is, depending on context, the three bit binary pattern `b111' can be taken
 to represent either `$7$' as an unsigned integer or `$-1$' as a signed, two's
 complement integer.

\subsubsection{Encoding Example}

The following example shows the state of an (8-bit) byte stream after several
 binary integers are encoded, including the location of the put pointer for the
 next bit to write to and the total length of the stream in bytes.

Encode the 4 bit unsigned integer value `12' (b1100) into an empty byte stream.

\begin{tabular}{rccccccccl}
                    &   &   &   &   &$\downarrow$&&&&            \\
        \hfill\vline& 7 & 6 & 5 & 4 & 3 & 2 & 1 & 0 &            \\\cline{1-9}
byte 0  \hfill\vline& 1 & 1 & 0 & 0 & 0 & 0 & 0 & 0 &$\leftarrow$\\
byte 1  \hfill\vline& 0 & 0 & 0 & 0 & 0 & 0 & 0 & 0 &            \\
byte 2  \hfill\vline& 0 & 0 & 0 & 0 & 0 & 0 & 0 & 0 &            \\
byte 3  \hfill\vline& 0 & 0 & 0 & 0 & 0 & 0 & 0 & 0 &            \\
\hfill$\vdots$\hfill\vline&\multicolumn{8}{c}{$\vdots$}&         \\
byte $n$\hfill\vline& 0 & 0 & 0 & 0 & 0 & 0 & 0 & 0 &
byte stream length: 1 byte
\end{tabular}
\vspace{\baselineskip}

Continue by encoding the 3 bit signed integer value `-1' (b111).

\begin{tabular}{rccccccccl}
                    &   &   &   &   &   &   &   &$\downarrow$&   \\
        \hfill\vline& 7 & 6 & 5 & 4 & 3 & 2 & 1 & 0 &            \\\cline{1-9}
byte 0  \hfill\vline& 1 & 1 & 0 & 0 & 1 & 1 & 1 & 0 &$\leftarrow$\\
byte 1  \hfill\vline& 0 & 0 & 0 & 0 & 0 & 0 & 0 & 0 &            \\
byte 2  \hfill\vline& 0 & 0 & 0 & 0 & 0 & 0 & 0 & 0 &            \\
byte 3  \hfill\vline& 0 & 0 & 0 & 0 & 0 & 0 & 0 & 0 &            \\
\hfill$\vdots$\hfill\vline&\multicolumn{8}{c}{$\vdots$}&         \\
byte $n$\hfill\vline& 0 & 0 & 0 & 0 & 0 & 0 & 0 & 0 &
byte stream length: 1 byte
\end{tabular}
\vspace{\baselineskip}

Continue by encoding the 7 bit integer value `17' (b0010001).

\begin{tabular}{rccccccccl}
                    &   &   &   &   &   &   &$\downarrow$&&      \\
        \hfill\vline& 7 & 6 & 5 & 4 & 3 & 2 & 1 & 0 &            \\\cline{1-9}
byte 0  \hfill\vline& 1 & 1 & 0 & 0 & 1 & 1 & 1 & 0 &            \\
byte 1  \hfill\vline& 0 & 1 & 0 & 0 & 0 & 1 & 0 & 0 &$\leftarrow$\\
byte 2  \hfill\vline& 0 & 0 & 0 & 0 & 0 & 0 & 0 & 0 &            \\
byte 3  \hfill\vline& 0 & 0 & 0 & 0 & 0 & 0 & 0 & 0 &            \\
\hfill$\vdots$\hfill\vline&\multicolumn{8}{c}{$\vdots$}&         \\
byte $n$\hfill\vline& 0 & 0 & 0 & 0 & 0 & 0 & 0 & 0 &
byte stream length: 2 bytes
\end{tabular}
\vspace{\baselineskip}

Continue by encoding the 13 bit integer value `6969' (b11011 00111001).

\begin{tabular}{rccccccccl}
                    &   &   &   &$\downarrow$&&&&   &            \\
        \hfill\vline& 7 & 6 & 5 & 4 & 3 & 2 & 1 & 0 &            \\\cline{1-9}
byte 0  \hfill\vline& 1 & 1 & 0 & 0 & 1 & 1 & 1 & 0 &            \\
byte 1  \hfill\vline& 0 & 1 & 0 & 0 & 0 & 1 & 1 & 1 &            \\
byte 2  \hfill\vline& 0 & 1 & 1 & 0 & 0 & 1 & 1 & 1 &            \\
byte 3  \hfill\vline& 0 & 0 & 1 & 0 & 0 & 0 & 0 & 0 &$\leftarrow$\\
\hfill$\vdots$\hfill\vline&\multicolumn{8}{c}{$\vdots$}&         \\
byte $n$\hfill\vline& 0 & 0 & 0 & 0 & 0 & 0 & 0 & 0 &
byte stream length: 4 bytes
\end{tabular}
\vspace{\baselineskip}

\subsubsection{Decoding Example}

The following example shows the state of the (8-bit) byte stream encoded in the
 previous example after several binary integers are decoded, including the
 location of the get pointer for the next bit to read.

Read a two bit unsigned integer from the example encoded above.

\begin{tabular}{rccccccccl}
                    &   &   &$\downarrow$&&&&   &   &            \\
        \hfill\vline& 7 & 6 & 5 & 4 & 3 & 2 & 1 & 0 &            \\\cline{1-9}
byte 0  \hfill\vline& 1 & 1 & 0 & 0 & 1 & 1 & 1 & 0 &$\leftarrow$\\
byte 1  \hfill\vline& 0 & 1 & 0 & 0 & 0 & 1 & 1 & 1 &            \\
byte 2  \hfill\vline& 0 & 1 & 1 & 0 & 0 & 1 & 1 & 1 &            \\
byte 3  \hfill\vline& 0 & 0 & 1 & 0 & 0 & 0 & 0 & 0 &
byte stream length: 4 bytes
\end{tabular}
\vspace{\baselineskip}

Value read: 3 (b11).

Read another two bit unsigned integer from the example encoded above.

\begin{tabular}{rccccccccl}
                    &   &   &   &   &$\downarrow$&&&            &            \\
        \hfill\vline& 7 & 6 & 5 & 4 & 3 & 2 & 1 & 0 &            \\\cline{1-9}
byte 0  \hfill\vline& 1 & 1 & 0 & 0 & 1 & 1 & 1 & 0 &$\leftarrow$\\
byte 1  \hfill\vline& 0 & 1 & 0 & 0 & 0 & 1 & 1 & 1 &            \\
byte 2  \hfill\vline& 0 & 1 & 1 & 0 & 0 & 1 & 1 & 1 &            \\
byte 3  \hfill\vline& 0 & 0 & 1 & 0 & 0 & 0 & 0 & 0 &
byte stream length: 4 bytes
\end{tabular}
\vspace{\baselineskip}

Value read: 0 (b00).

Two things are worth noting here.
\begin{itemize}
\item
Although these four bits were originally written as a single four-bit integer,
 reading some other combination of bit-widths from the bitstream is well
 defined.
No artificial alignment boundaries are maintained in the bitstream.
\item
The first value is the integer `$3$' only because the context stated we were
 reading an unsigned integer.
Had the context stated we were reading a signed integer, the returned value
 would have been the integer `$-1$'.
\end{itemize}

\subsubsection{End-of-Packet Alignment}

The typical use of bitpacking is to produce many independent byte-aligned
 packets which are embedded into a larger byte-aligned container structure,
 such as an Ogg transport bitstream.
Externally, each bitstream encoded as a byte stream must begin and end on a
 byte boundary.
Often, the encoded bitstream is not an integer number of bytes, and so there is
 unused space in the last byte of a packet.

Unused space in the last byte of a packet is always zeroed during the encoding
 process.
Thus, should this unused space be read, it will return binary zeroes.
There is no marker pattern or stuffing bits that will allow the decoder to
 obtain the exact size, in bits, of the original bitstream.
This knowledge is not required for decoding.

Attempting to read past the end of an encoded packet results in an
 `end-of-packet' condition.
Any further read operations after an `end-of-packet' condition shall also
 return `end-of-packet'.
Unlike Vorbis, Theora does not use truncated packets as a normal mode of
 operation.
Therefore if a decoder encounters the `end-of-packet' condition during normal
 decoding, it may attempt to use the bits that were read to recover as much of
 encoded data as possible, signal a warning or error, or both.

\subsubsection{Reading Zero Bit Integers}

Reading a zero bit integer returns the value `$0$' and does not increment
 the stream pointer.
Reading to the end of the packet, but not past the end, so that an
 `end-of-packet' condition is not triggered, and then reading a zero bit
 integer shall succeed, returning `$0$', and not trigger an `end-of-packet'
 condition.
Reading a zero bit integer after a previous read sets the `end-of-packet'
 condition shall fail, also returning `end-of-packet'.

\bibliography{spec}

\end{document}
